% Limitations and threats to validity.
\section{Limitations and threats to validity}\label{sec:limitations}

\subsection{Corpus and queries}
\uapar{} imitates a personal archive with public material: Wikipedia articles stand in for notes and rendered articles for office documents, so the text is more uniform and better edited than real notes, and nothing in it is truly personal. The 200 queries were written by one model under written rules, not by people looking for their own documents, and every answerable target is a note. The separate 50-query extension targets documents, tables and repository files (Section~\ref{sec:results-extension}), but is again model-written. It contains 49 distinct wordings, so the 50 records are not fully independent; a deterministic unique-wording sensitivity is reported separately. The no-answer queries have no lexical match by construction, which decides the lexical systems' result on them in advance. Topics follow the author's interests; a different archive could favour different systems.

\subsection{Model-generated relevance labels}
Text relevance labels are generated by language models rather than an independent human panel. Agreement, target-document checks and deliberately unanswerable queries provide diagnostic evidence but are not a complete gold standard. Gemini and Gemma may share errors because they share a provider. Reviewing disagreements with Gemini cannot remove this dependence; it also exposes the reviewer to the initial grades. We therefore preserve the original labels and report the sensitivity of retrieval metrics and system ordering to the judge and to the review. The primary judge was selected after examining judge diagnostics, a post hoc methodological choice. This edition deliberately includes no human relevance audit. Its text-retrieval conclusions therefore remain conditional on the model judgements; independent human validation is future work.

\subsection{Query and candidate coverage}
Although the corpus is heterogeneous, all authored answerable targets in the original 200-query set are notes. The separate extension adds a small sample of documents, tables and repository files. Neither set establishes general quality for arbitrary target formats, images or recordings. The E5 baseline searches a smaller text-only document universe, whereas the EG2 indexes include media. The control on identical candidates (Table~\ref{tab:encoder-control}) changes EG2's nDCG@10 by $-0.0015$, a small point effect in this control. This does not remove all candidate-coverage uncertainty: the control runs were not pooled, and 0.29\% of their top-10 documents are unjudged.

Only the union of system top-10 results is fully judged. Relevant documents outside that pool may exist. Recall at 100 is therefore relative to the known relevant pool; treating unjudged documents as nonrelevant can favour pooled systems and does not estimate exhaustive corpus recall. A failure to reject a performance difference does not establish equivalence.

\subsection{Single-author field study}
The field study (Section~\ref{sec:field}) uses one person's archive and query sets that the author wrote during development of the system, partly to test it; they are not a random sample of real information needs, and the system was tuned on the same archive. Its numbers show behaviour in deployment, not generalizable effectiveness.

\subsection{Measurements}
All costs come from one laptop in an ordinary desktop session, with other programs running; the load before each step is logged but not controlled. The query timer includes embedding and chunk search but excludes document collapse, path resolution and interface rendering. Process-level resource counters span the whole measured pass, including collapse, so their scope differs from the per-query timer. CPU indexing was measured on a fifth of the GPU sample. The scaling experiment uses synthetic vectors derived from the real ones, and HNSW recall on them need not equal recall on a real collection of that size.

\subsection{Construct validity}
Judges see the title and the first 1\,200 characters of a document, not the whole document, so a document relevant only further down can be graded too low; the same snippet is shown for every system, so the bias does not favour one system directly but favours documents whose relevance shows early. Grades 0/1/2 compress a judgement that people make differently: a model's ``partly relevant'' is not a person's. The primary metric, nDCG@10 over pooled labels, rewards systems for ranking judged-relevant documents high; it does not measure whether a person finds what they remember faster, which is what an archive search is for.

\subsection{Source-pair matching tasks}
Image and speech tasks identify the paired source item or sentence. Those pairings do not exhaustively judge every semantically related candidate. Captions and parallel transcripts also differ from independently paraphrased archive queries. The reported point estimates have no computed confidence intervals.
